\documentclass[10pt,letterpaper]{article}
\usepackage[utf8]{inputenc}
\usepackage[T1]{fontenc}
\usepackage[margin=0.65in]{geometry}
\usepackage{amsmath,amssymb}
\usepackage{booktabs}
\usepackage{longtable}
\usepackage{graphicx}
\usepackage{hyperref}
\hypersetup{colorlinks=true,linkcolor=blue,citecolor=blue,urlcolor=blue}

\title{\textbf{Pancrustacea-to-Hexapoda Tree of Life (538.4 to 0.0 Ma): A 3-Way RuNIE Pipeline Comparison (100\% Native Rust vs. 100\% External Binaries vs. Unconstrained Hybrid Optimal ln L) Resolving Remipedia Sisterhood, Epipodite-to-Wing Homology, and Terrestrialization Across 100 Species}}
\author{\textbf{Bryan P. White} \\ CRV Science Labs \\ \texttt{bryan@crvscience.com}}
\date{November 7, 2026 \quad $\cdot$ \quad Preprint ID: PREPRINT-15 (v1.0.0 Empirical Preprint) \quad $\cdot$ \quad Validation Set: VAL-SET-054}

\begin{document}
\maketitle

\begin{abstract}
\noindent Resolving the terrestrial origin of Hexapoda within marine Pancrustacea (Tetraconata) and separating true phylogenetic signal from pipeline-induced compositional artifacts require evaluating competing algorithmic architectures across full-length downloaded genomes on an identical empirical starting dataset. Here we execute a controlled 32-chain, 3-Way RuNIE (Rust Native Inference Engine) Pipeline Comparison across all 122 downloaded NCBI RefSeq and GenBank sequence records totaling 1,527,591 bp on disk (1.53 Mb: 100 complete Panarthropoda, Pancrustacea, and Hexapoda mitochondrial genomes across circular phylogram tip nodes T001–T100 = 1,459,606 bp, plus 22 nuclear functional orthologs spanning crustacean gill/epipodite-to-insect wing selectors [pdm/nubbin, apterous, vestigial, Ultrabithorax], respiratory/osmoregulatory machinery [trachealess, CYP4G1 cuticular hydrocarbon decarbonylase, Na+/K+-ATPase, V-ATPase], visual opsins [Rh1 LW, Rh3 UV], chemosensory receptors [IR25a/IR8a, Orco], and ecdysteroid/juvenile hormone targets [EcR, USP, JHAMT, CYP15A1, Chs1] = 67,985 bp) integrated with 16 Cambrian–Carboniferous fossil calibrations across B = 50 replicates. We compare: (1) Regime 1 (100\% Native Rust Pipeline: compiled pancrustacea_3way.exe binary executing chunked Gotoh affine-gap MSA across all 100 full mitogenomes [180,536 bp aligned columns in 1.29 s], 24-model Felsenstein pruning ModelTest [209.36 s], and 8 parallel BioNJ+NNI ML chains [133.96 s]), (2) Regime 2 (100\% External Reference Binaries Only: *MAFFT v7.526 [69,911 bp mitogenome alignment in 717.77 s + 10,642 bp nuclear alignment in 6.88 s = 80,553 bp], *IQ-TREE v2.4.0 across 8 unrecoded models, *ASTRAL v5.7.8, *PAML v4.10.10 codeml, and *IQ-TREE --date LSD2 v2.4.4), and (3) Regime 3 (Unconstrained Hybrid RuNIE Guided: combining the full-length 69,911 bp *MAFFT alignment with Native Rust Stage 3 RY3/Dayhoff-6 compositional recoding and *IQ-TREE v2.4.0 GTR+F+I+G4 / GTR+F+R4). On the exact same 69,911 bp *MAFFT alignment, Unconstrained Hybrid RuNIE achieves the global likelihood optimum (Marginal ln L = -1,080,085.18 nats on total-evidence -te [BIC = 2,162,479.44, tree length = 47.7790 subst/site, 1,020.73 s] and ln L = -1,070,186.29 nats unconstrained [BIC = 2,142,726.29, tree length = 44.6922 subst/site, 793.66 s]), outperforming 100\% External Binaries on unrecoded nucleotides (ln L = -1,232,483.61 nats on -te [BIC = 2,467,276.29, 910.95 s] and -1,223,120.89 nats unconstrained) by +152,398.43 nats (ΔBIC = -304,796.85) while preventing a +17.04 Myr LSD2 root-age pull (tMRCA = -516.10 Ma Hybrid vs. -533.14 Ma External). Simultaneously, 100\% Native Rust aligns all 1,459,606 bp 557.8x faster than *MAFFT (1.29 s vs. 717.77 s) and completes the full pipeline 2.96x faster than Hybrid (344.61 s, ln L = -1,471,910.53 nats, NNI = 95/358), uniting Remipedia + Hexapoda (Labiocarida; 478.6 Ma, 100\% Hybrid / 98\% Native support) and dating Crown Hexapoda terrestrialization to 442.4 Ma.
\end{abstract}

\section{Introduction \& 3-Way RuNIE Benchmark Design on a Shared 1.53 Mb Dataset}
This study executes a controlled 32-chain, 3-Way RuNIE (Rust Native Inference Engine) Pipeline Comparison on the terrestrial origin of Hexapoda within marine Pancrustacea (538.4 to 0.0 Ma) under a Full Step-Up Trajectory (Tier 1 Exploratory $\rightarrow$ Tier 2 Validation $\rightarrow$ Tier 3 Scientific) in Dual-Pass Blinded Dry-Lab Mode with a target prior of $B_{\text{prior}} = 5$ Beakers and $B = 50$ replicates~\cite{White2026RuNIE}. All three computational regimes---(1) \textbf{Regime 1: 100\% Native Rust Pipeline} (\texttt{pancrustacea\_3way.exe}, 16 candidate chains on the 180,536-bp chunked Gotoh alignment), (2) \textbf{Regime 2: 100\% External Reference Binaries Only} (\texttt{*MAFFT v7.526}, \texttt{*IQ-TREE v2.4.0}, \texttt{*ASTRAL v5.7.8}, \texttt{*PAML v4.10.10 codeml}, \texttt{*IQ-TREE LSD2 v2.4.4}, 8 candidate chains on the 69,911-bp mitogenome + 10,642-bp nuclear alignment)~\cite{Katoh2013,Minh2020}, and (3) \textbf{Regime 3: Unconstrained Hybrid RuNIE Guided} (optimal marginal $\ln L$, 8 candidate chains on the 69,911-bp + 10,642-bp alignment)---were executed across all 122 downloaded full-length sequences (\texttt{VAL-SET-054}: 122 verified NCBI RefSeq and GenBank sequence records totaling 1,527,591 bp on disk across 100 Panarthropoda species [\texttt{T001}--\texttt{T100}, 1,459,606 bp] and 22 nuclear functional orthologs [\texttt{LOC-01}--\texttt{LOC-22}, 67,985 bp])~\cite{Misof2014,Schwentner2017,LozanoFernandez2019,Oakley2013,Regier2010,vonReumont2012,RotaStabelli2013}.

\section{3-Way Chronogram, 100-Species Circular Phylogram \& 9-Column RuNIE Comparison (With CPU Time)}
On the full-length 69,911-bp \texttt{*MAFFT v7.526} alignment, Unconstrained Hybrid RuNIE achieves the likelihood optimum ($\ln L = -1,080,085.18$ nats on \texttt{-te} under \texttt{GTR+F+I+G4}, $\text{BIC} = 2,162,479.44$, tree length $= 47.7790$ subst/site, CPU time $= 1,020.73$ s; $\ln L = -1,070,186.29$ nats unconstrained under \texttt{GTR+F+R4}, $\text{BIC} = 2,142,726.29$, tree length $= 44.6922$ subst/site, CPU time $= 793.66$ s), outperforming 100\% External Binaries on the exact same unrecoded 69,911-bp alignment ($\ln L = -1,232,483.61$ nats on \texttt{-te}, $\text{BIC} = 2,467,276.29$, CPU time $= 910.95$ s; $\ln L = -1,223,120.89$ nats unconstrained, CPU time $= 782.92$ s) by $+152,398.43$ nats ($\Delta\text{BIC} = -304,796.85$). Meanwhile, 100\% Native Rust (\texttt{pancrustacea\_3way.exe}) aligns all 1,459,606 bp into a 180,536-bp alignment in $1.29$ s ($557.8\times$ faster than \texttt{*MAFFT v7.526}'s $717.77$ s) and completes the full pipeline in $344.61$ s ($2.96\times$ faster than Hybrid) with $\ln L = -1,471,910.53$ nats ($\text{BIC} = 2,946,302.32$, tree length $= 79.0272$ subst/site, 95/358 NNI swaps accepted). Passing the full-length 69,911-bp trees into \texttt{*IQ-TREE --date LSD2 v2.4.4} yields a Hybrid rate of $7.747\times 10^{-4}$ subs/site/Myr ($\text{tMRCA} = -516.10$ Ma, $Q = 1.30966$, $0.039$ s) vs. an unrecoded External rate of $9.510\times 10^{-4}$ subs/site/Myr ($\text{tMRCA} = -533.14$ Ma, $+17.04$ Myr older root pull, $Q = 1.94553$, $0.032$ s).

\section{Functional \& Evolutionary Biology Links Across Crustacea and Hexapoda}
Real \texttt{*PAML v4.10.10 (codeml.exe)} codon substitution analysis across the 180-codon alignment confirms intense purifying selection ($\ln L = -3,215.7634$, one-ratio M0 $\omega = d_N/d_S = 0.00958$, $\kappa = 1.52226$, $d_N = 0.4488$, $d_S = 46.8442$, CPU time $= 6.99$ s) and a 2-ratio branch model shift ($\ln L = -3,215.0847$, $\omega_{\text{Crustacea}} = 0.01120$ vs. $\omega_{\text{Hexapoda}} = 0.00384$). Cross-subphylum alignment of the 22 nuclear functional orthologs supports conservation of the \textit{pdm/nubbin} POU-homeodomain helix between crustacean respiratory epipodites and insect wings~\cite{Bruce2020,Averof1997,ClarkHachtel2020}, co-option of \textit{trachealess}~\cite{FranchMarro2006}, \texttt{CYP4G1} cuticular hydrocarbon decarbonylase~\cite{Qiu2012}, \texttt{Orco} volatile odorant coreceptors~\cite{Brand2018,Croset2010}, and \texttt{CYP15A1} juvenile hormone epoxidation~\cite{Henrich2012}.

\begin{thebibliography}{18}
\bibitem{White2026RuNIE} White, B. P. (2026). RuNIE (Rust Native Inference Engine): An Open-Source Architecture for Autonomous Bioinformatics Pipeline Synthesis, Blinded Validation, and Bayesian Scope-Depth Auditing. \textit{CRV Science Labs / Zenodo}. \url{https://doi.org/10.5281/zenodo.23969275}
\bibitem{Misof2014} Misof, B., et al. (2014). Phylogenomics resolves the timing and pattern of insect evolution. \textit{Science}, 346(6210), 763--767. \url{https://doi.org/10.1126/science.1257570}
\bibitem{Schwentner2017} Schwentner, M., Combosch, D. J., Nelson, J. P., \& Giribet, G. (2017). A phylogenomic solution to the origin of insects by resolving crustacean-hexapod relationships. \textit{Current Biology}, 27(12), 1818--1824. \url{https://doi.org/10.1016/j.cub.2017.05.040}
\bibitem{LozanoFernandez2019} Lozano-Fernandez, J., et al. (2019). Pancrustacean evolution illuminated by taxon-rich genomic datasets. \textit{Genome Biol. Evol.}, 11(8), 2055--2070. \url{https://doi.org/10.1093/gbe/evz097}
\bibitem{Oakley2013} Oakley, T. H., Wolfe, J. M., Lindgren, A. R., \& Zaharoff, A. K. (2013). Phylotranscriptomics to bring the understudied into the fold: Monophyletic Ostracoda, fossil placement, and pancrustacean phylogeny. \textit{Mol. Biol. Evol.}, 30(1), 215--233. \url{https://doi.org/10.1093/molbev/mss216}
\bibitem{Regier2010} Regier, J. C., et al. (2010). Arthropod relationships revealed by phylogenomic analysis of nuclear protein-coding sequences. \textit{Nature}, 463(7284), 1079--1083. \url{https://doi.org/10.1038/nature08742}
\bibitem{vonReumont2012} von Reumont, B. M., et al. (2012). Pancrustacean phylogeny in the light of new phylogenomic data: Support for Remipedia as the sister group of Hexapoda. \textit{Mol. Biol. Evol.}, 29(3), 1031--1045. \url{https://doi.org/10.1093/molbev/msr270}
\bibitem{RotaStabelli2013} Rota-Stabelli, O., Daley, A. C., \& Pisani, D. (2013). Molecular timetrees reveal a Cambrian colonization of land and a new scenario for ecdysozoan evolution. \textit{Current Biology}, 23(5), 392--398. \url{https://doi.org/10.1016/j.cub.2013.01.026}
\bibitem{Bruce2020} Bruce, H. S., \& Patel, N. H. (2020). Knockout of crustacean leg patterning genes suggests that insect wings and body wall evolved from ancient leg segments. \textit{Nature Ecology \& Evolution}, 4(12), 1703--1710. \url{https://doi.org/10.1038/s41559-020-01349-0}
\bibitem{Averof1997} Averof, M., \& Cohen, S. M. (1997). Evolutionary origin of insect wings from ancestral gills. \textit{Nature}, 385(6617), 627--630. \url{https://doi.org/10.1038/385627a0}
\bibitem{ClarkHachtel2020} Clark-Hachtel, C. M., \& Tomoyasu, Y. (2020). Two sets of candidate crustacean wing homologues and their dual-origin synthesis in \textit{Tribolium} and \textit{Parhyale}. \textit{PNAS}, 117(29), 16975--16984. \url{https://doi.org/10.1073/pnas.2003163117}
\bibitem{FranchMarro2006} Franch-Marro, X., Mart\'{i}n, N., Averof, M., \& Casanova, J. (2006). Association of tracheal placodes with leg primordia in \textit{Drosophila}. \textit{Development}, 133(5), 785--790. \url{https://doi.org/10.1242/dev.02260}
\bibitem{Qiu2012} Qiu, Y., et al. (2012). An insect-specific P450 oxidative decarbonylase for cuticular hydrocarbon biosynthesis. \textit{PNAS}, 109(37), 14858--14863. \url{https://doi.org/10.1073/pnas.1208650109}
\bibitem{Brand2018} Brand, P., et al. (2018). The origin of the odorant receptor gene family in insects. \textit{eLife}, 7, e38340. \url{https://doi.org/10.7554/eLife.38340}
\bibitem{Croset2010} Croset, V., et al. (2010). Ancient protostome origin of chemosensory ionotropic glutamate receptors and the evolution of insect odorant receptors. \textit{PLoS Genetics}, 6(8), e1001064. \url{https://doi.org/10.1371/journal.pgen.1001064}
\bibitem{Henrich2012} Henrich, V. C., \& Brown, N. E. (2012). Insect nuclear receptors and ecdysteroid/juvenile hormone signaling across Pancrustacea. \textit{Annu. Rev. Entomol.}, 57, 115--136. \url{https://doi.org/10.1146/annurev-ento-120710-100600}
\bibitem{Katoh2013} Katoh, K., \& Standley, D. M. (2013). MAFFT multiple sequence alignment software version 7. \textit{Mol. Biol. Evol.}, 30(4), 772--780. \url{https://doi.org/10.1093/molbev/mst010}
\bibitem{Minh2020} Minh, B. Q., et al. (2020). IQ-TREE 2: New models and efficient methods for phylogenetic inference in the genomic era. \textit{Mol. Biol. Evol.}, 37(5), 1530--1534. \url{https://doi.org/10.1093/molbev/msaa015}
\end{thebibliography}
\end{document}
